% =====================================================================
% Capítulo 1 — Esboço de seção/interlúdio
% Sugestão de inserção: imediatamente após o parágrafo que termina com
% "[...] fazer alianças \parencite{stengers2010cosmopolitics1}." (tese
% como inscrição tecnocientífica), e antes do parágrafo que abre com
% "De certo modo, em sintonia com as críticas de Haraway e Stengers [...]".
%
% Observação técnica: para o build do LaTeX encontrar a figura, copie
%   infranodus/infranodus_cap1_focus.png
% para
%   figuras/cap.1/infranodus_cap1_focus.png
% (mantendo a convenção do capítulo). O caminho abaixo já assume essa
% localização. Se preferir manter a figura em `infranodus/`, basta ajustar
% o argumento de \includegraphics.
% =====================================================================

\subsection*{Interlúdio: a tese como inscrição que se mede a si mesma}

Se esta tese é uma inscrição tecnocientífica, então ela mesma pode ser
submetida ao tipo de tratamento que ela aplica aos seus interlocutores:
pode ser convertida em traço, em diagrama, em superfície combinável.
\textcite{Latour1986Visualization} sugere que a força de uma inscrição
está em ser \textit{móvel imutável}, isto é, transportável e
recombinável em uma única superfície plana onde inscrições heterogêneas
passam a ser comparadas. Os centros de cálculo, segue Latour, agem à
distância porque acumulam essas superfícies. O que ofereço aqui é um
gesto experimental e propositadamente recursivo: tomo o texto deste
capítulo como objeto e o submeto a uma análise de rede textual no
estilo do \textit{InfraNodus} \parencite{paranyushkin2019infranodus},
em que o texto é representado como grafo de co-ocorrência de termos e
as comunidades desse grafo são lidas como tópicos latentes.\footnote{A
análise foi conduzida com um \textit{script} em Python que reproduz a
metodologia descrita por \textcite{paranyushkin2019infranodus}: o texto
do capítulo, depurado dos comandos \LaTeX{} e com as notas de rodapé
reincorporadas ao corpus, foi tokenizado em \textit{lemas} e percorrido
por uma janela deslizante de quatro \textit{tokens}; cada par de termos
co-ocorrentes na janela recebeu um peso decrescente com a distância
(3, 2, 1). O grafo ponderado foi podado para os 180 termos mais
frequentes e para arestas com peso $\geq 2$; em seguida, foram
detectadas comunidades pelo algoritmo de Louvain ponderado e calculadas
medidas de \textit{degree} ponderado e \textit{betweenness centrality}.
Os artefatos da análise estão depositados em
\texttt{infranodus/} no repositório desta tese (\texttt{.gexf} para
\textit{Gephi}, \texttt{.png}, \texttt{.csv}, e relatório
\texttt{.md}). Por se tratar de uma reprodução do método e não do
\textit{software} comercial, registro o gesto como \textit{análise
no estilo InfraNodus}.}

A leitura desse grafo é, ela própria, uma operação de tradução: cada
nó é um termo, cada aresta é uma co-ocorrência, e cada cor é uma
comunidade detectada por um algoritmo cuja lógica é tão alheia à
hermenêutica antropológica quanto as redes neurais que descrevo nos
Capítulos~\ref{capitulo3} e~\ref{capitulo4}. E, ainda assim, ele me
devolve algo. Os termos com maior \textit{degree} ponderado neste
capítulo são, em ordem, \textit{pesquisa}, \textit{ciência},
\textit{rede}, \textit{pesquisador(a)}, \textit{modelo},
\textit{humano}, \textit{objeto}, \textit{inteligência},
\textit{artificial}, \textit{campo} e \textit{etnografia}. Os termos
que mais costuram comunidades distintas (maior \textit{betweenness}) são
\textit{ciência}, \textit{pesquisa}, \textit{rede} e \textit{pesquisador(a)}:
não termos teóricos densos, mas justamente os termos genéricos que
permitem o trânsito entre o vocabulário ator-rede, o vocabulário das
ciências sociais e o vocabulário da inteligência artificial. A própria
condição de \textit{ponte}, no sentido que \textcite{Latour1986Visualization}
dá ao termo, parece coincidir aqui com a condição de \textit{tradutor}
no sentido latouriano mais amplo \parencite{Latour2011CienciaEmAcao}.

A detecção de comunidades estabilizou seis tópicos
latentes:\footnote{Os rótulos abaixo correspondem aos quatro a seis
termos de maior \textit{degree} dentro de cada comunidade; preservo a
forma lematizada produzida pelo \textit{script} para evitar tradução
interpretativa precoce.}
(i)~\textit{escrita} e \textit{composição} em torno de
\textit{pesquisador}, \textit{modelo}, \textit{objeto}, \textit{Claude};
(ii)~vocabulário \textbf{ator-rede} em torno de \textit{rede},
\textit{Latour}, \textit{agência}, \textit{ator}, \textit{diagrama},
\textit{distribuída};
(iii)~\textbf{ciências sociais e campo} em torno de \textit{ciência},
\textit{sociais}, \textit{dado}, \textit{conhecimento},
\textit{cientista};
(iv)~\textbf{IA e etnografia} em torno de \textit{pesquisa},
\textit{inteligência}, \textit{artificial}, \textit{etnografia},
\textit{construção}, \textit{técnico};
(v)~\textbf{humano-máquina-método} em torno de \textit{humano},
\textit{método}, \textit{existência}, \textit{máquina},
\textit{parcial};
(vi)~um cluster periférico em torno de \textit{Simondon},
\textit{técnica} e \textit{alienação}.

\begin{figure}[!htb]
\centering
\includegraphics[width=1.0\textwidth,keepaspectratio]{figuras/cap.1/infranodus_cap1_focus.png}
\caption[Rede textual do Capítulo~\ref{capitulo1} (análise no estilo InfraNodus)]{Rede
textual do Capítulo~\ref{capitulo1} (análise no estilo
\textit{InfraNodus}). Os nós são os 100 termos com maior
\textit{degree} ponderado; o tamanho do nó é proporcional ao
\textit{degree}; a cor indica a comunidade detectada por Louvain
ponderado; arestas com peso~$\geq 3$. Versão de alta resolução,
arquivo \texttt{.gexf} para o \textit{Gephi} e relatório completo
disponíveis em \texttt{infranodus/} no repositório desta tese.}
\label{fig:infranodus-cap1}
\end{figure}

O que essa cartografia me dá não é uma confirmação, mas um conjunto de
\textit{chamadas de atenção}. A primeira é uma confirmação modesta:
os termos centrais do grafo coincidem com os conceitos que mobilizo
explicitamente, e os pontos de maior intermediação são justamente os
termos mais genéricos do capítulo --- o que sugere que a costura
argumentativa está sendo feita por palavras-coringa que ainda esperam
serem teorizadas com mais força. A segunda é uma surpresa: as duas
comunidades com a ligação ponderada mais fraca entre si são a das
\textbf{ciências sociais} (Tópico~iii) e a do par
\textbf{humano-máquina-método} (Tópico~v).
Em outras palavras, o grafo registra uma fissura textual entre o
vocabulário das ciências sociais com que escrevo e o vocabulário da
simetria humano/não-humano com que prometo escrever. A formulação de
\textcite{Law2004After} sobre o \textit{method assemblage} ajuda a
nomear o que isso quer dizer: a \textit{otherness} apagada para que a
presença se sustente é, neste capítulo, a articulação entre os dois
registros que ainda preciso fazer trabalhar juntos. A terceira chamada
diz respeito à comunidade periférica em torno de Simondon: ela está no
grafo, mas mal conectada, sinalizando uma promessa teórica que o
capítulo abre e que os capítulos subsequentes precisam costurar à rede
maior, sob risco de virar nota de rodapé errante.

Reconheço o paradoxo do gesto. Aplicar análise de rede a um capítulo
que invoca a teoria ator-rede é arriscar tomar a metáfora pelo método;
e a similaridade entre os dois usos da palavra \textit{rede} é
superficial: a rede ator-rede é uma cadeia de associações
heterogêneas e ontologicamente plana, ao passo que a rede textual aqui
gerada é um grafo bidimensional de co-ocorrências de \textit{tokens}
ponderado por proximidade.\footnote{A confusão entre os dois sentidos
de \textit{rede} é discutida por \textcite[pp.~129--132]{latour2005reassembling},
que distingue \textit{rede} como ferramenta de inscrição e como
diagnóstico ontológico. Mantenho a distinção: o que o
\textit{InfraNodus} produz é uma inscrição; o que esta tese descreve
é uma rede sociotécnica. As duas se encontram apenas no gesto de
inscrever uma na outra.} Mas talvez seja exatamente isso que justifica
o gesto: a rede textual não substitui a etnografia, ela funciona como
mais uma inscrição que esta tese mobiliza, no mesmo registro das
imagens, dos diagramas e do site que a acompanham. Como qualquer
inscrição, ela ordena, hierarquiza, inclui e exclui --- ou, em outras
palavras, faz política ontológica \parencite[p.~143]{Law2004After}. E
talvez essa seja a versão mais honesta do gesto reflexivo que esta
tese tenta sustentar: medir o próprio texto pelos mesmos instrumentos
com que se compromete a medir o C4AI, sabendo que medir é também
produzir aquilo que se mede.

% =====================================================================
% Sugestão de entrada bibliográfica (caso ainda não esteja no .bib):
%
% @article{paranyushkin2019infranodus,
%   author    = {Paranyushkin, Dmitry},
%   title     = {{InfraNodus}: Generating Insight Using Text Network Analysis},
%   booktitle = {The World Wide Web Conference (WWW '19)},
%   year      = {2019},
%   pages     = {3584--3589},
%   publisher = {ACM},
%   doi       = {10.1145/3308558.3314123}
% }
% =====================================================================
